% Part: incompleteness
% Chapter: incompleteness-provability
% Section: lob-thm

\documentclass[../../../include/open-logic-section]{subfiles}

\begin{document}

\olfileid[hi]{inc}{inp}{lob}

\olsection{लोब (L\"ob) का प्रमेय}

किसी सिद्धांत~$\Th{T}$ का ग्योडेल (G\"odel) वाक्य
$\lnot \OProv[\Th{T}](y)$ का कोई नियत-बिंदु है; अर्थात् ऐसा
!!a{sentence}~$!G$ कि
\[
\Th{T} \Proves \lnot \OProv[\Th{T}](\gn{!G}) \liff !G.
\]
वह !!{derivable} नहीं है, क्योंकि यदि $\Th{T} \Proves !G$, तो (क)
!!{derivability} प्रतिबंध~(1) से
$\Th{T} \Proves \OProv[\Th{T}](\gn{!G})$, और (ख)
$\Th{T} \Proves !G$ को
$\Th{T} \Proves \lnot \OProv[\Th{T}](\gn{!G}) \liff !G$ के साथ मिलाने
पर $\Th{T} \Proves \lnot \OProv[\Th{T}](\gn{!G})$ मिलता है; अतः
$\Th{T}$ असंगत होगा। अब $\OProv[\Th{T}](y)$ के किसी नियत-बिंदु की स्थिति
के बारे में पूछना स्वाभाविक है; अर्थात् ऐसे !!a{sentence}~$!H$ के बारे में
कि
\[
\Th{T} \Proves \OProv[\Th{T}](\gn{!H}) \liff !H.
\]
यदि वह !!{derivable} होता, तो प्रतिबंध~(1) से
$\Th{T} \Proves \OProv[\Th{T}](\gn{!H})$; किंतु ऊपर की समतुल्यता पर मोडस
पोनेन्स लगाने से भी यही निष्कर्ष मिलता है। अतः कम-से-कम ग्योडेल (G\"odel) वाक्य वाले
तर्क से यह निष्कर्ष नहीं मिलता कि $\Th{T}$ असंगत है। निस्संदेह इससे यह
नहीं दिखता कि $\Th{T}$,~$!H$ को !!{derive} \emph{करता} है।

प्रश्न को थोड़ा व्यापक बनाकर हम आगे बढ़ सकते हैं। नियत-बिंदु समतुल्यता की
बाएँ-से-दाएँ दिशा $\OProv[\Th{T}](\gn{!H}) \lif !H$ एक सामान्य योजना का
उदाहरण है, जिसे \emph{प्रतिबिंबन सिद्धांत} कहते हैं:
$\OProv[\Th{T}](\gn{!A}) \lif !A$। इसका यह नाम इसलिए है कि एक अर्थ में यह
व्यक्त करता है कि $\Th{T}$ अपनी !!{derive} कर सकने वाली बातों पर
``प्रतिबिंबित'' कर सकता है; मूलतः वह किसी भी~$!A$ के लिए कहता है, ``यदि
$\Th{T}$,~$!A$ को !!{derive} कर सकता है, तो~$!A$ सत्य है।'' निस्संदेह यह
केवल सुदृढ़ सिद्धांतों के लिए सत्य है, और इससे संकेत मिलता है कि सिद्धांत
सामान्यतः इसके प्रत्येक उदाहरण को !!{derive} नहीं करेंगे। तो कोई
(पर्याप्त प्रबल और !!{derivability} प्रतिबंधों को संतुष्ट करने वाला)
सिद्धांत कौन-से उदाहरण !!{derive} कर सकता है? निश्चिततः वे सभी जिनमें
$!A$ स्वयं !!{derivable} है। और बस वही, जैसा अगला परिणाम दिखाता है।

\begin{thm}\ollabel{thm:lob}
$\Th{T}$,~$\Th{Q}$ का कोई !!a{axiomatizable} विस्तार हो, और मान लें
$\OProv[\Th{T}](y)$ कोई ऐसा सूत्र है जो \olref[2in]{sec} के प्रतिबंध
P1--P3 को संतुष्ट करता है। यदि $\Th{T}$,
$\OProv[\Th{T}](\gn{!A}) \lif !A$ को !!{derive} करता है, तो वास्तव में
$\Th{T}$,~$!A$ को !!{derive} करता है।
\end{thm}

भिन्न शब्दों में, यदि $\Th{T} \Proves/ !A$, तो
$\Th{T} \Proves/ \OProv[\Th{T}](\gn{!A}) \lif !A$। यह परिणाम लोब (L\"ob)
 का प्रमेय कहलाता है।

\begin{explain}
लोब (L\"ob) के प्रमेय के प्रमाण की मार्गदर्शक कल्पना सांता क्लॉज़ के
अस्तित्व का एक चतुर प्रमाण है। (यदि आपको वह निष्कर्ष पसंद न हो, तो अपनी
पसंद का कोई भी दूसरा निष्कर्ष रख सकते हैं।) प्रमाण यह है:
\begin{enumerate}
\item $X$ वह वाक्य हो, ``यदि $X$ सत्य है, तो सांता क्लॉज़ का अस्तित्व है।''
\item मान लें $X$ सत्य है।
\item तब उसकी कही हुई बात लागू होती है; अर्थात् हमारे पास है: यदि $X$ सत्य
  है, तो सांता क्लॉज़ का अस्तित्व है।
\item चूँकि हम $X$ को सत्य मान रहे हैं, (2) और~(3) से मोडस पोनेन्स द्वारा
  निष्कर्ष निकाल सकते हैं कि सांता क्लॉज़ का अस्तित्व है।
\item हम मान्यता~(2), ``$X$ सत्य है,'' से (4), ``सांता क्लॉज़ का अस्तित्व
  है,'' व्युत्पन्न करने में सफल हुए हैं। सप्रतिबंध प्रमाण से हमने दिखाया:
  ``यदि $X$ सत्य है, तो सांता क्लॉज़ का अस्तित्व है।''
\item किंतु यह केवल वाक्य~$X$ है। अतः हमने दिखाया कि $X$ सत्य है।
\item किंतु तब ऊपर के तर्क (2)--(4) से सांता क्लॉज़ का अस्तित्व है।
\end{enumerate}
इस विचार को औपचारिक बनाकर, ``सत्य है'' के स्थान पर ``!!{derivable} है'' और
``सांता क्लॉज़ का अस्तित्व है'' के स्थान पर~$!A$ रखने से लोब (L\"ob) के
प्रमेय का प्रमाण मिलता है। युक्ति है नियत-बिंदु उपपत्ति को
!!{formula}~$\OProv[\Th{T}](y) \lif !A$ पर लागू करना। उसका नियत-बिंदु ऊपर
की रूपरेखा के वाक्य~$X$ के संगत है।
\end{explain}

\begin{proof}[\olref{thm:lob} का प्रमाण]
मान लें $!A$ ऐसा !!a{sentence} है कि $\Th{T}$,
$\OProv[\Th{T}](\gn{!A}) \lif !A$ को !!{derive} करता है। $!B(y)$ को
!!{formula}~$\OProv[\Th{T}](y) \lif !A$ मानें और नियत-बिंदु उपपत्ति का
उपयोग करके कोई !!a{sentence}~$!D$ ऐसा पाएँ कि $\Th{T}$,
$!D \liff !B(\gn{!D})$ को !!{derive} करता है। तब निम्न में से प्रत्येक
$\Th{T}$ में !!{derivable} है:
\begin{align}
  & !D \liff (\OProv[\Th{T}](\gn{!D}) \lif !A) \ollabel{L-1}\\
  & \qquad \text{$!D$,~$!B(y)$ का नियत-बिंदु है}\notag \\
  & !D \lif (\OProv[\Th{T}](\gn{!D}) \lif !A) \ollabel{L-2}\\
  & \qquad\text{\olref{L-1} से}\notag\\
  & \OProv[\Th{T}](\gn{!D \lif (\OProv[\Th{T}](\gn{!D}) \lif !A)}) \ollabel{L-3}\\
  & \qquad \text{\olref{L-2} से, प्रतिबंध P1 द्वारा}\notag \\
  & \OProv[\Th{T}](\gn{!D}) \lif \OProv[\Th{T}](\gn{\OProv[\Th{T}](\gn{!D}) \lif !A})
  \ollabel{L-4}\\
  &\qquad \text{\olref{L-3} से, प्रतिबंध P2 द्वारा}\notag \\
  & \OProv[\Th{T}](\gn{!D}) \lif (\OProv[\Th{T}](\gn{\OProv[\Th{T}](\gn{!D})}) \lif \OProv[\Th{T}](\gn{!A})) \ollabel{L-5}\\
  &\qquad \text{\olref{L-4} से, फिर P2 का उपयोग करके} \notag\\
& \OProv[\Th{T}](\gn{!D}) \lif \OProv[\Th{T}](\gn{\OProv[\Th{T}](\gn{!D})}) \ollabel{L-6}\\
  & \qquad\text{!!{derivability} प्रतिबंध P3 से} \notag\\
  & \OProv[\Th{T}](\gn{!D}) \lif \OProv[\Th{T}](\gn{!A}) \ollabel{L-7} \\
  &\qquad\text{\olref{L-5} और \olref{L-6} से}\notag\\
  & \OProv[\Th{T}](\gn{!A}) \lif !A \ollabel{L-8}\\
  &\qquad\text{प्रमेय की मान्यता से} \notag\\
  & \OProv[\Th{T}](\gn{!D}) \lif !A \ollabel{L-9}\\
  &\qquad\text{\olref{L-7} और \olref{L-8} से}\notag\\
  & (\OProv[\Th{T}](\gn{!D}) \lif !A) \lif !D \ollabel{L-10}\\
  & \qquad \text{\olref{L-1} से}\notag \\
  & !D \ollabel{L-11}\\
  & \qquad\text{\olref{L-9} और \olref{L-10} से}\notag \\
  & \OProv[\Th{T}](\gn{!D}) \ollabel{L-12}\\
  & \qquad\text{\olref{L-11} से, प्रतिबंध~P1 द्वारा}\notag \\
  & !A \qquad\qquad\text{\olref{L-8} और \olref{L-12} से।}\notag
\end{align}
\end{proof}

लोब (L\"ob) का प्रमेय उपलब्ध होने पर द्वितीय अपूर्णता-प्रमेय का छोटा प्रमाण
मिलता है (उन सिद्धांतों के लिए जिनमें प्रतिबंध P1--P3 को संतुष्ट करने वाला
!!a{derivability} विधेय है): यदि
$\Th{T} \Proves \OProv[\Th{T}](\gn{\lfalse}) \lif \lfalse$, तो
$\Th{T} \Proves \lfalse$। यदि $\Th{T}$ संगत है, तो
$\Th{T} \Proves/ \lfalse$। अतः
$\Th{T} \Proves/ \OProv[\Th{T}](\gn{\lfalse}) \lif \lfalse$, अर्थात्
$\Th{T} \Proves/ \OCon[\Th{T}]$। हम इसे यह दिखाने के लिए भी लागू कर
सकते हैं कि $\OProv[\Th{T}](x)$ का नियत-बिंदु~$!H$,~!!{derivable} है।
क्योंकि
\begin{align*}
  \Th{T} & \Proves \OProv[\Th{T}](\gn{!H}) \liff !H\\
  \intertext{और विशेषतः}
    \Th{T} & \Proves \OProv[\Th{T}](\gn{!H}) \lif !H,
\end{align*}
इसलिए लोब (L\"ob) के प्रमेय से $\Th{T} \Proves !H$।

% Going in the other direction, for homework I
% may ask you to work through a short proof of L\"ob's theorem, using
% the second incompleteness theorem instead of the fixed-point lemma.

\begin{prob}
$\Th{T}$ कोई संगणनीयतः अभिगृहीतीकृत सिद्धांत हो और
$\OProv[\Th{T}]$,~$\Th{T}$ का कोई !!a{derivability} विधेय हो। निम्न चार
कथनों पर विचार करें:
\begin{enumerate}
\item यदि $T \Proves !A$, तो $T \Proves \OProv[\Th{T}](\gn{!A})$।
\item $T \Proves !A \lif \OProv[\Th{T}](\gn{!A})$।
\item यदि $T \Proves \OProv[\Th{T}](\gn{!A})$, तो $T \Proves !A$।
\item $T \Proves \OProv[\Th{T}](\gn{!A}) \lif !A$।
\end{enumerate}
इनमें से प्रत्येक कथन किन प्रतिबंधों में सत्य है?
\end{prob}

\end{document}
